\documentclass[12pt,a4paper]{article}
\usepackage[utf8]{inputenc}
\usepackage[spanish]{babel}
\usepackage{geometry}
\usepackage{listings}
\usepackage{xcolor}
\usepackage{parskip}
\usepackage{amsmath}
\usepackage{hyperref}

\geometry{margin=2.5cm}

\lstset{
  basicstyle=\ttfamily\small,
  keywordstyle=\color{blue}\bfseries,
  commentstyle=\color{gray},
  breaklines=true, frame=single, language=Java,
  numbers=left, numberstyle=\tiny\color{gray}
}

\title{\textbf{Algoritmo de Wang --- Demostrador de Teoremas v1}\\
  \large Documentación Técnica}
\author{Inteligencia Artificial}
\date{}

\begin{document}
\maketitle
\tableofcontents
\newpage

\section{Descripción del Proyecto}
Primera versión del demostrador de teoremas basado en el \textbf{Algoritmo de Wang} (1960). Determina automáticamente si un argumento lógico proposicional es una tautología usando el cálculo de secuentes.

\section{Fundamento Teórico}

\subsection{Secuentes}
Un \textbf{secuente} tiene la forma $\Gamma \vdash \Delta$ donde:
\begin{itemize}
  \item $\Gamma$ (antecedente): conjunto de premisas
  \item $\Delta$ (consecuente): conjunto de conclusiones
\end{itemize}
El secuente es \textbf{válido} si, cuando todas las fórmulas de $\Gamma$ son verdaderas, al menos una de $\Delta$ también lo es.

\subsection{Axioma de cierre}
Una rama se \textbf{cierra} (es válida) cuando la misma fórmula aparece en el antecedente y en el consecuente: $A, \Gamma \vdash \Delta, A$.

\section{Reglas de Inferencia}

\begin{center}
\begin{tabular}{|l|l|l|}
\hline
\textbf{Regla} & \textbf{Nombre} & \textbf{Transformación} \\
\hline
$\neg L$ & Negación-izq. & $\neg A, \Gamma \vdash \Delta \Rightarrow \Gamma \vdash A, \Delta$ \\
$\neg R$ & Negación-der. & $\Gamma \vdash \neg A, \Delta \Rightarrow A, \Gamma \vdash \Delta$ \\
$\wedge L$ & Conjunción-izq. & $A \wedge B, \Gamma \vdash \Delta \Rightarrow A, B, \Gamma \vdash \Delta$ \\
$\rightarrow L$ & Implicación-izq. & $A \rightarrow B, \Gamma \vdash \Delta \Rightarrow \neg A, B, \Gamma \vdash \Delta$ \\
$\rightarrow R$ & Implicación-der. & $\Gamma \vdash A \rightarrow B \Rightarrow A, \Gamma \vdash B$ \\
\hline
\end{tabular}
\end{center}

\section{Implementación del Algoritmo}

\begin{lstlisting}
boolean esValido(List<String> antecedente,
                 List<String> consecuente) {
    // Caso base: axioma de cierre
    for (String f : antecedente)
        if (consecuente.contains(f)) return true;

    // Aplicar reglas al antecedente
    for (String f : antecedente) {
        if (f.startsWith("!")) { // negacion izq
            List<String> nuevoAnt = sin(antecedente, f);
            List<String> nuevoCons = con(consecuente,
                                          f.substring(1));
            return esValido(nuevoAnt, nuevoCons);
        }
        if (f.contains("->")) { // implicacion izq
            String[] partes = f.split("->");
            List<String> nuevoAnt = sin(antecedente, f);
            nuevoAnt.add("!" + partes[0]);
            nuevoAnt.add(partes[1]);
            return esValido(nuevoAnt, consecuente);
        }
        // ... más reglas
    }
    return false; // no se pudo cerrar
}
\end{lstlisting}

\section{Parser de Fórmulas}

El parser convierte la entrada del usuario en un árbol sintáctico:
\begin{itemize}
  \item \texttt{!A} → negación $\neg A$
  \item \texttt{A\&B} → conjunción $A \wedge B$
  \item \texttt{A|B} → disyunción $A \vee B$
  \item \texttt{A->B} → implicación $A \rightarrow B$
\end{itemize}

\section{Interfaz Gráfica}

\begin{itemize}
  \item Campos de texto para ingresar premisas y conclusión.
  \item Panel de salida que muestra cada paso de la demostración.
  \item Resultado final: \textbf{VÁLIDO} (tautología) o \textbf{NO VÁLIDO} (refutable).
\end{itemize}

\end{document}
